\section*{Заключение}
\addcontentsline{toc}{section}{Заключение}

В ходе выполнения курсового проекта было разработано веб-приложение для ведения семейного бюджета \texttt{MoneyManage}. Приложение позволяет семьям совместно контролировать доходы и расходы, устанавливать бюджеты для участников, загружать чеки и получать наглядную статистику.

В результате выполнения работы были решены все поставленные задачи:

\begin{enumerate}
    \item Проведён анализ существующих программных продуктов (Дзен-мани, CoinKeeper, Money Lover), выявивший их недостатки: платные премиум-функции, слабая поддержка совместной работы, отсутствие полноценных веб-версий. Обоснован выбор технологического стека: Python + Flask, PostgreSQL, Bootstrap 5, Chart.js.
    
    \item Определены требования целевой аудитории (семьи с совместным бюджетом) и спроектирована функциональность приложения, включая диаграммы вариантов использования и User Story.
    
    \item Разработана модель данных, включающая 8 сущностей (\texttt{Family}, \texttt{User}, \texttt{Category}, \texttt{Transaction}, \texttt{Receipt}, \texttt{CategoryLimit}, \texttt{UserBudget}, \texttt{DashboardStats}), построена ER-диаграмма.
    
    \item Созданы макеты страниц (wireframes) для всех категорий пользователей: администратора, владельца семьи и участника.
    
    \item Реализована система аутентификации и авторизации с тремя ролями (\texttt{admin}, \texttt{owner}, \texttt{member}) с использованием \texttt{Flask-Login} и \texttt{bcrypt}. При регистрации владельца автоматически создаются категории по умолчанию.
    
    \item Разработан CRUD-интерфейс для транзакций, категорий и участников. Реализованы фильтрация транзакций по дате и поиску, установка лимитов на категории и персональных бюджетов.
    
    \item Реализованы механизмы загрузки чеков (изображения, PDF) с сохранением на сервере и хранением пути в БД. Добавлен экспорт отчётов в формате CSV.
    
    \item Реализована система визуализации: круговая диаграмма расходов по категориям и столбчатая диаграмма динамики баланса (Chart.js). Для оптимизации производительности используется кэширование ежедневных балансов в модели \texttt{DashboardStats}.
    
    \item Проведено тестирование приложения, код размещён в репозитории GitHub: \url{https://github.com/userownmaa/web_coursework_Momina_241-327}. Оформлена отчётная документация.
\end{enumerate}

Разработанное приложение представляет собой полноценный программный продукт, готовый к развёртыванию на хостинге. Дальнейшее развитие проекта может включать: добавление мобильной версии, интеграцию с банковскими API для автоматического импорта транзакций, расширенную аналитику с прогнозированием расходов, а также улучшение системы уведомлений о превышении лимитов.
